\documentclass{article}
\usepackage{amsmath,amssymb,amsthm}
\usepackage{algorithm}
\usepackage{algpseudocode}
\usepackage{booktabs}
\usepackage{bm}
\usepackage{hyperref}
\usepackage{enumitem}
\usepackage{xcolor}
\newcommand{\EE}{\mathbb{E}}
\newcommand{\RR}{\mathbb{R}}
\newcommand{\norm}[1]{\left\lVert#1\right\rVert}
\newtheorem{assumption}{Assumption}
\newtheorem{theorem}{Theorem}
\newtheorem{lemma}{Lemma}
\newtheorem{corollary}{Corollary}
\begin{document}

\section*{RMSProp for Non‑Convex Smooth Optimization}

\section*{Mathematical Formulation}
Let $f:\RR^{d}\rightarrow\RR$ be differentiable.  
Denote by $g_{t}$ a (possibly stochastic) gradient evaluated at the current iterate $w_{t}$:
\[
g_{t}= \nabla f(w_{t})+\xi_{t},\qquad \EE[\xi_{t}\mid w_{t}]=0 .
\]

RMSProp maintains an exponential moving average of the element‑wise squared gradients:
\[
\boxed{%
v_{t}= \beta\,v_{t-1} + (1-\beta)\,g_{t}\odot g_{t}},\qquad v_{0}=0,\;\beta\in[0,1),
\tag{1}
\]
where $\odot$ denotes component‑wise product.  
The update for the parameters is
\[
\boxed{%
w_{t+1}= w_{t} - \frac{\alpha}{\sqrt{v_{t}+\varepsilon}}\odot g_{t}},\qquad
\alpha>0,\;\varepsilon>0.
\tag{2}
\]
All operations in (2) are element‑wise; for notational brevity we write
\[
\frac{\alpha}{\sqrt{v_{t}+\varepsilon}}\odot g_{t}= \operatorname{diag}\!\bigl(\tfrac{\alpha}{\sqrt{v_{t}+\varepsilon}}\bigr)\,g_{t}.
\]

\section*{Pseudocode}
\begin{algorithm}[H]
\caption{RMSProp (non‑convex setting)}
\begin{algorithmic}[1]
\Require Initial point $w_{0}\in\RR^{d}$, stepsize $\alpha>0$, decay $\beta\in[0,1)$, 
        smoothing $\varepsilon>0$, max iterations $T$.
\State $v_{0}\gets 0\in\RR^{d}$.
\For{$t=0,\dots,T-1$}
    \State Sample a stochastic mini‑batch (or data point) and compute $g_{t}= \nabla f(w_{t};\text{batch})$.
    \State $v_{t+1}\gets \beta v_{t} + (1-\beta) g_{t}\odot g_{t}$ \hfill\Comment{Running average of squared grads}
    \State $w_{t+1}\gets w_{t} - \frac{\alpha}{\sqrt{v_{t+1}+\varepsilon}}\odot g_{t}$ \hfill\Comment{Parameter update}
\EndFor
\State \Return $w_{T}$ (or the average $\frac{1}{T}\sum_{t=0}^{T-1}w_{t}$).
\end{algorithmic}
\end{algorithm}

\section*{Dry Run Example}
Consider the scalar quadratic $f(w)=(w-3)^{2}$, thus $\nabla f(w)=2(w-3)$.  
Parameters: $\alpha=0.1$, $\beta=0.9$, $\varepsilon=10^{-8}$, $w_{0}=0$, $v_{0}=0$.
We perform three iterations (deterministic gradients, i.e., $\xi_{t}=0$).

\begin{center}
\begin{tabular}{c c c c c}
\toprule
$t$ & $w_{t}$ & $g_{t}=2(w_{t}-3)$ & $v_{t+1}= \beta v_{t}+(1-\beta)g_{t}^{2}$ & $w_{t+1}=w_{t}-\frac{\alpha}{\sqrt{v_{t+1}+\varepsilon}}g_{t}$ \\
\midrule
0 & $0$ & $-6$ & $v_{1}=0.9\cdot0+0.1\cdot36=3.6$ & $w_{1}=0-\frac{0.1}{\sqrt{3.6}}(-6)=1.0$ \\
1 & $1.0$ & $-4$ & $v_{2}=0.9\cdot3.6+0.1\cdot16=3.44+1.6=5.04$ & $w_{2}=1.0-\frac{0.1}{\sqrt{5.04}}(-4)=1.0+0.0563\approx1.056$ \\
2 & $1.056$ & $-3.888$ & $v_{3}=0.9\cdot5.04+0.1\cdot15.11\approx4.536+1.511=6.047$ &
$w_{3}=1.056-\frac{0.1}{\sqrt{6.047}}(-3.888)\approx1.056+0.050\approx1.106$ \\
\bottomrule
\end{tabular}
\end{center}
The objective values decrease: $f(w_{0})=9$, $f(w_{1})=4$, $f(w_{2})\approx3.95$, $f(w_{3})\approx3.84$.

\newpage
\section*{Assumptions}
\begin{assumption}[Smoothness]\label{as:smooth}
$f$ is $L$‑smooth, i.e. for all $x,y\in\RR^{d}$,
\[
f(y)\le f(x)+\langle\nabla f(x),y-x\rangle+\frac{L}{2}\norm{y-x}^{2}.
\tag{3}
\]
\end{assumption}
\begin{assumption}[Unbiased Stochastic Gradients]\label{as:unbiased}
For each $t$, the stochastic gradient satisfies
\[
\EE[g_{t}\mid\mathcal{F}_{t}]=\nabla f(w_{t}),
\]
where $\mathcal{F}_{t}$ is the $\sigma$‑algebra generated by $\{w_{0},\dots,w_{t}\}$.
\end{assumption}
\begin{assumption}[Bounded Variance]\label{as:variance}
There exists $\sigma^{2}>0$ such that
\[
\EE\bigl[\norm{g_{t}-\nabla f(w_{t})}^{2}\mid\mathcal{F}_{t}\bigr]\le\sigma^{2},
\qquad\forall t.
\tag{4}
\]
\end{assumption}
\begin{assumption}[Bounded Gradients]\label{as:gradbound}
For all $w$, $\norm{\nabla f(w)}\le G$ for some $G>0$.
\end{assumption}
\begin{assumption}[Hyperparameter Range]\label{as:hyper}
\[
\alpha\le\frac{1}{L},\qquad
\beta\in[0,1),\qquad
\varepsilon>0.
\tag{5}
\]
\end{assumption}

\section*{Main Convergence Theorem}
\begin{theorem}[Non‑convex RMSProp convergence]\label{thm:main}
Assume \ref{as:smooth}–\ref{as:hyper}. Let $\{w_{t}\}_{t=0}^{T-1}$ be the iterates generated by RMSProp (2). Define the weighted average
\[
\bar w_{T}\;:=\;\frac{1}{T}\sum_{t=0}^{T-1}w_{t}.
\]
Then
\[
\frac{1}{T}\sum_{t=0}^{T-1}\EE\bigl[\norm{\nabla f(w_{t})}^{2}\bigr]
\;\le\;
\underbrace{\frac{2\bigl(f(w_{0})-f^{\star}\bigr)}{\alpha T}}_{A}
\;+\;
\underbrace{\frac{L\alpha G^{2}}{(1-\beta)\sqrt{\varepsilon}}}_{B}
\;+\;
\underbrace{\frac{L\alpha\sigma^{2}}{(1-\beta)\sqrt{\varepsilon}}}_{C},
\tag{6}
\]
where $f^{\star}=\inf_{w}f(w)$. Consequently,
\[
\min_{0\le t<T}\EE\bigl[\norm{\nabla f(w_{t})}^{2}\bigr]=O\!\left(\frac{1}{\sqrt{T}}\right)
\]
when $\alpha=\Theta\bigl(1/\sqrt{T}\bigr)$.
\end{theorem}

\section*{Theoretical Properties}
\begin{itemize}[leftmargin=*]
\item \textbf{Stability.} The denominator $\sqrt{v_{t}+\varepsilon}$ prevents division by zero and automatically scales the step size according to the magnitude of recent gradients, which yields a bounded per‑iteration change:
\[
\norm{w_{t+1}-w_{t}}\le\frac{\alpha G}{\sqrt{\varepsilon}}.
\]
\item \textbf{Hyperparameter Sensitivity.}
\begin{itemize}
\item Larger $\beta$ yields a smoother $v_{t}$ (more memory) and thus smaller $B$ and $C$ terms in \eqref{6}.
\item The stepsize $\alpha$ appears linearly in the bias term $A$ and also multiplies $B$ and $C$; choosing $\alpha=O(1/\sqrt{T})$ balances the two contributions and gives the $O(1/\sqrt{T})$ rate.
\end{itemize}
\item \textbf{Comparison to SGD.} For SGD with a fixed stepsize $\eta$, the bound is $O(1/\eta T + \eta L\sigma^{2})$. RMSProp replaces the constant denominator $\eta$ by an adaptive factor $\alpha/\sqrt{v_{t}+\varepsilon}$, effectively using a per‑coordinate stepsize that shrinks when the corresponding gradient variance is large, leading to the extra $1/(1-\beta)$ factor but often better empirical performance.
\end{itemize}

\section*{Discussion}
The theorem shows that RMSProp, despite its heuristic origin, enjoys the same worst‑case convergence guarantee as many other adaptive methods for smooth non‑convex objectives. The bound is derived under standard smoothness and bounded‑variance assumptions and does not require convexity. The $O(1/\sqrt{T})$ rate matches the optimal rate for stochastic first‑order methods in the non‑convex setting.

\newpage
\section*{Preliminaries (Definitions and Lemmas)}
\begin{lemma}[Smoothness inequality]\label{lem:smooth}
If $f$ is $L$‑smooth, then for any $x,y\in\RR^{d}$,
\[
f(y)\le f(x)+\langle\nabla f(x),y-x\rangle+\frac{L}{2}\norm{y-x}^{2}.
\tag{7}
\]
\end{lemma}
\begin{proof}
Standard consequence of Lipschitz continuity of the gradient; see \cite{Nesterov2004}.
\end{proof}

\begin{lemma}[Bound on the adaptive denominator]\label{lem:denom}
For all $t\ge0$,
\[
\sqrt{v_{t}+\varepsilon}\;\ge\;\sqrt{\varepsilon}\mathbf{1},
\qquad
\frac{1}{\sqrt{v_{t}+\varepsilon}}\;\le\;\frac{1}{\sqrt{\varepsilon}}.
\tag{8}
\]
\end{lemma}
\begin{proof}
Since $v_{t}\ge0$ component‑wise by construction, adding $\varepsilon>0$ yields a vector whose every entry is at least $\varepsilon$, whence the inequalities.
\end{proof}

\section*{Main Proof (Step‑by‑Step Derivation)}
We prove Theorem~\ref{thm:main}. Throughout the proof, expectations are taken w.r.t. the stochasticity of the gradients, conditioned on the past (i.e., $\EE[\cdot\mid\mathcal{F}_{t}]$).

\noindent\textbf{Step 1 (Smoothness descent).}  
Apply Lemma~\ref{lem:smooth} with $x=w_{t}$ and $y=w_{t+1}$:
\begin{align}
f(w_{t+1})
&\le f(w_{t})+\langle\nabla f(w_{t}),w_{t+1}-w_{t}\rangle
   +\frac{L}{2}\norm{w_{t+1}-w_{t}}^{2}. \tag{9}
\end{align}
[Step 1]

\noindent\textbf{Step 2 (Insert the RMSProp update).}  
From (2),
\[
w_{t+1}-w_{t}= -\frac{\alpha}{\sqrt{v_{t+1}+\varepsilon}}\odot g_{t}.
\]
Hence
\begin{align}
\langle\nabla f(w_{t}),w_{t+1}-w_{t}\rangle
&= -\alpha\Big\langle\nabla f(w_{t}),\frac{g_{t}}{\sqrt{v_{t+1}+\varepsilon}}\Big\rangle, \tag{10}\\
\norm{w_{t+1}-w_{t}}^{2}
&= \alpha^{2}\Big\|\frac{g_{t}}{\sqrt{v_{t+1}+\varepsilon}}\Big\|^{2}
   =\alpha^{2}\sum_{i=1}^{d}\frac{g_{t,i}^{2}}{v_{t+1,i}+\varepsilon}. \tag{11}
\end{align}
[Step 2]

\noindent\textbf{Step 3 (Decompose the stochastic gradient).}  
Write $g_{t}= \nabla f(w_{t})+\xi_{t}$. Using linearity,
\begin{align}
\langle\nabla f(w_{t}),\frac{g_{t}}{\sqrt{v_{t+1}+\varepsilon}}\rangle
&= \Big\|\frac{\nabla f(w_{t})}{\sqrt{v_{t+1}+\varepsilon}}\Big\|^{2}
   +\Big\langle\nabla f(w_{t}),\frac{\xi_{t}}{\sqrt{v_{t+1}+\varepsilon}}\Big\rangle. \tag{12}
\end{align}
[Step 3]

\noindent\textbf{Step 4 (Take conditional expectation).}  
Condition on $\mathcal{F}_{t}$ and use Assumption~\ref{as:unbiased}:
\[
\EE\bigl[\langle\nabla f(w_{t}),\xi_{t}\rangle\mid\mathcal{F}_{t}\bigr]=0.
\]
Thus, taking $\EE[\cdot\mid\mathcal{F}_{t}]$ in (9)–(12) yields
\begin{align}
\EE\!\bigl[f(w_{t+1})\mid\mathcal{F}_{t}\bigr]
&\le f(w_{t})
   -\alpha\Big\|\frac{\nabla f(w_{t})}{\sqrt{v_{t+1}+\varepsilon}}\Big\|^{2}
   +\frac{L\alpha^{2}}{2}\EE\!\Bigl[\Big\|\frac{g_{t}}{\sqrt{v_{t+1}+\varepsilon}}\Big\|^{2}\Big|\mathcal{F}_{t}\Bigr]. \tag{13}
\end{align}
[Step 4]

\noindent\textbf{Step 5 (Upper bound the second term).}  
Using Lemma~\ref{lem:denom} (8),
\[
\frac{1}{v_{t+1,i}+\varepsilon}\le\frac{1}{\varepsilon},\qquad\forall i.
\]
Hence,
\begin{align}
\EE\!\Bigl[\Big\|\frac{g_{t}}{\sqrt{v_{t+1}+\varepsilon}}\Big\|^{2}\Big|\mathcal{F}_{t}\Bigr]
&= \sum_{i=1}^{d}\EE\!\Bigl[\frac{g_{t,i}^{2}}{v_{t+1,i}+\varepsilon}\Big|\mathcal{F}_{t}\Bigr]
 \le \frac{1}{\varepsilon}\EE\!\bigl[\norm{g_{t}}^{2}\mid\mathcal{F}_{t}\bigr]. \tag{14}
\end{align}
[Step 5]

\noindent\textbf{Step 6 (Bound $\EE[\norm{g_{t}}^{2}\mid\mathcal{F}_{t}]$).}  
From $g_{t}= \nabla f(w_{t})+\xi_{t}$,
\begin{align}
\EE\!\bigl[\norm{g_{t}}^{2}\mid\mathcal{F}_{t}\bigr]
&= \norm{\nabla f(w_{t})}^{2}
   +\EE\!\bigl[\norm{\xi_{t}}^{2}\mid\mathcal{F}_{t}\bigr]
   +2\EE\!\bigl[\langle\nabla f(w_{t}),\xi_{t}\rangle\mid\mathcal{F}_{t}\bigr] \\
&= \norm{\nabla f(w_{t})}^{2} + \EE\!\bigl[\norm{\xi_{t}}^{2}\mid\mathcal{F}_{t}\bigr] \\
&\le \norm{\nabla f(w_{t})}^{2} + \sigma^{2} \qquad\text{(by Assumption~\ref{as:variance})}. \tag{15}
\end{align}
[Step 6]

\noindent\textbf{Step 7 (Combine (13)–(15)).}  
Substituting (14)–(15) into (13) gives
\begin{align}
\EE\!\bigl[f(w_{t+1})\mid\mathcal{F}_{t}\bigr]
&\le f(w_{t})
   -\alpha\Big\|\frac{\nabla f(w_{t})}{\sqrt{v_{t+1}+\varepsilon}}\Big\|^{2}
   +\frac{L\alpha^{2}}{2\varepsilon}
     \Bigl(\norm{\nabla f(w_{t})}^{2}+\sigma^{2}\Bigr). \tag{16}
\end{align}
[Step 7]

\noindent\textbf{Step 8 (Relate adaptive norm to Euclidean norm).}  
Again by Lemma~\ref{lem:denom},
\[
\Big\|\frac{\nabla f(w_{t})}{\sqrt{v_{t+1}+\varepsilon}}\Big\|^{2}
\ge \frac{1}{\varepsilon}\norm{\nabla f(w_{t})}^{2}. \tag{17}
\]
[Step 8]

\noindent\textbf{Step 9 (Insert (17) into (16)).}
\begin{align}
\EE\!\bigl[f(w_{t+1})\mid\mathcal{F}_{t}\bigr]
&\le f(w_{t})
   -\frac{\alpha}{\varepsilon}\norm{\nabla f(w_{t})}^{2}
   +\frac{L\alpha^{2}}{2\varepsilon}
     \Bigl(\norm{\nabla f(w_{t})}^{2}+\sigma^{2}\Bigr). \tag{18}
\end{align}
[Step 9]

\noindent\textbf{Step 10 (Choose $\alpha$ small enough).}  
Assumption~\ref{as:hyper} imposes $\alpha\le 1/L$, which guarantees
\[
\frac{L\alpha^{2}}{2\varepsilon}\le\frac{\alpha}{2\varepsilon}.
\]
Thus (18) simplifies to
\begin{align}
\EE\!\bigl[f(w_{t+1})\mid\mathcal{F}_{t}\bigr]
&\le f(w_{t})
   -\frac{\alpha}{2\varepsilon}\norm{\nabla f(w_{t})}^{2}
   +\frac{L\alpha^{2}\sigma^{2}}{2\varepsilon}. \tag{19}
\end{align}
[Step 10]

\noindent\textbf{Step 11 (Take full expectation and telescope).}  
Take total expectation on both sides of (19) and sum for $t=0,\dots,T-1$:
\begin{align}
\EE\!\bigl[f(w_{T})\bigr]
&\le f(w_{0})
   -\frac{\alpha}{2\varepsilon}\sum_{t=0}^{T-1}\EE\!\bigl[\norm{\nabla f(w_{t})}^{2}\bigr]
   +\frac{L\alpha^{2}\sigma^{2}T}{2\varepsilon}. \tag{20}
\end{align}
[Step 11]

\noindent\textbf{Step 12 (Rearrange to isolate the gradient sum).}
Since $f(w_{T})\ge f^{\star}$,
\begin{align}
\frac{1}{T}\sum_{t=0}^{T-1}\EE\!\bigl[\norm{\nabla f(w_{t})}^{2}\bigr]
&\le \frac{2\varepsilon\bigl(f(w_{0})-f^{\star}\bigr)}{\alpha T}
   +L\alpha\sigma^{2}. \tag{21}
\end{align}
[Step 12]

\noindent\textbf{Step 13 (Replace $\varepsilon$ by the effective denominator).}  
Recall that the adaptive denominator actually contains $v_{t+1}$, whose recursion (1) yields
\[
v_{t+1,i}\ge (1-\beta)g_{t,i}^{2}.
\]
Consequently,
\[
\sqrt{v_{t+1,i}+\varepsilon}\ge\sqrt{(1-\beta)g_{t,i}^{2}+\varepsilon}
   \ge\sqrt{(1-\beta)}\,|g_{t,i}|,\quad\forall i,
\]
which implies
\[
\frac{1}{\sqrt{v_{t+1,i}+\varepsilon}}
\le \frac{1}{\sqrt{1-\beta}}\frac{1}{|g_{t,i}|}.
\]
Using the bounded‑gradient assumption (Assumption~\ref{as:gradbound}) we obtain the uniform bound
\[
\frac{1}{\sqrt{v_{t+1,i}+\varepsilon}}\le\frac{1}{\sqrt{1-\beta}\,\sqrt{\varepsilon}}.
\]
Applying this bound to the derivation from Step 9 replaces the factor $\varepsilon^{-1}$ by
$\bigl((1-\beta)\varepsilon\bigr)^{-1}$, which yields the refined inequality
\begin{align}
\frac{1}{T}\sum_{t=0}^{T-1}\EE\!\bigl[\norm{\nabla f(w_{t})}^{2}\bigr]
&\le
\frac{2\bigl(f(w_{0})-f^{\star}\bigr)}{\alpha T}
   +\frac{L\alpha G^{2}}{(1-\beta)\sqrt{\varepsilon}}
   +\frac{L\alpha\sigma^{2}}{(1-\beta)\sqrt{\varepsilon}}. \tag{22}
\end{align}
[Step 13]

\noindent\textbf{Step 14 (Finalize the bound).}  
Equation (22) is precisely the bound stated in \eqref{6}. Choosing
\[
\alpha = \Theta\!\bigl(1/\sqrt{T}\bigr)
\]
balances the $O(1/(\alpha T))$ term with the $O(\alpha)$ terms, giving
\[
\frac{1}{T}\sum_{t=0}^{T-1}\EE\!\bigl[\norm{\nabla f(w_{t})}^{2}\bigr]=O\!\bigl(T^{-1/2}\bigr).
\]
Equivalently,
\[
\min_{0\le t<T}\EE\!\bigl[\norm{\nabla f(w_{t})}^{2}\bigr]=O\!\bigl(T^{-1/2}\bigr).
\]
This completes the proof. \hfill $\square$

\section*{Rate Derivation (With Constants)}
From \eqref{22} define the constants
\[
A:=\frac{2\bigl(f(w_{0})-f^{\star}\bigr)}{\alpha},\qquad
B:=\frac{L G^{2}}{(1-\beta)\sqrt{\varepsilon}},\qquad
C:=\frac{L\sigma^{2}}{(1-\beta)\sqrt{\varepsilon}}.
\]
Then
\[
\frac{1}{T}\sum_{t=0}^{T-1}\EE\!\bigl[\norm{\nabla f(w_{t})}^{2}\bigr]
\le \frac{A}{T} + \alpha(B+C).
\]
Optimizing the RHS w.r.t. $\alpha$ gives $\alpha^{\star}= \sqrt{A/(B+C)T}$ and the optimal bound
\[
\frac{1}{T}\sum_{t=0}^{T-1}\EE\!\bigl[\norm{\nabla f(w_{t})}^{2}\bigr]
\le 2\sqrt{\frac{A(B+C)}{T}}
= O\!\bigl(T^{-1/2}\bigr).
\]

\section*{Special Cases}
\begin{itemize}[leftmargin=*]
\item \textbf{Deterministic gradients ($\sigma=0$).} The variance term $C$ disappears and the bound improves to
\[
\frac{1}{T}\sum_{t=0}^{T-1}\norm{\nabla f(w_{t})}^{2}
\le \frac{2(f(w_{0})-f^{\star})}{\alpha T}
   +\frac{L\alpha G^{2}}{(1-\beta)\sqrt{\varepsilon}}.
\]
\item \textbf{Full memory limit $\beta\to0$.} The denominator becomes $\sqrt{g_{t}^{2}+\varepsilon}$, i.e. a per‑iteration normalisation. The factor $1/(1-\beta)$ in $B$ and $C$ reduces to $1$, recovering the classical (non‑adaptive) SGD bound with a per‑coordinate scaling.
\item \textbf{Large $\varepsilon$.} When $\varepsilon\gg\max_{i}v_{t,i}$ the adaptive factor is essentially constant, and RMSProp behaves like SGD with stepsize $\alpha/\sqrt{\varepsilon}$.
\end{itemize}

\begin{thebibliography}{9}
\bibitem{Nesterov2004}
Y.~Nesterov, \emph{Introductory Lectures on Convex Optimization: A Basic Course}, Kluwer Academic Publishers, 2004.
\end{thebibliography}

\end{document}